% ============================================================================
%   RÉSUMÉ (Français) — Aivancity: 250-300 mots, mots-clés, autonome
% ============================================================================

\chapter*{Résumé}
\addcontentsline{toc}{chapter}{Résumé}
\markboth{Résumé}{}

\selectlanguage{french}

L'appariement patient--essai clinique est un goulot d'étranglement de 
la recherche médicale, et les systèmes récents fondés sur les grands 
modèles de langage (LLM) tels que TrialGPT visent à automatiser 
l'évaluation de l'éligibilité des patients face aux critères 
d'inclusion et d'exclusion. Au-delà de la précision brute, la 
fiabilité probabiliste de ces systèmes n'a cependant pas fait l'objet 
d'un audit systématique. Une revue de portée récente portant sur 
vingt-quatre systèmes LLM d'appariement n'en identifie aucun qui 
rapporte des métriques de calibration. Cette omission est lourde de 
conséquences : un modèle sur-confiant qui se trompe avec certitude 
est cliniquement pire qu'un modèle moins précis mais honnêtement 
incertain.

Nous avons mené le premier audit systématique de la calibration des 
LLM open source dans ce domaine. Douze configurations (six modèles 
croisés avec deux variantes de prompt) ont été évaluées sur les 
1\,015 paires patient--critère de TrialGPT-Criterion-Annotations. 
Quatre hypothèses pré-enregistrées ont été testées par tests de 
permutation avec correction de Bonferroni, et la stabilité a été 
vérifiée par validation croisée 5-fold et comparaison à des baselines 
naïves. Les distributions de probabilité sont extraites par scoring 
direct des log-probabilités plutôt que par génération libre.

Les douze configurations présentent un ECE au-dessus du seuil 
clinique de $0{,}05$, avec des valeurs entre $0{,}124$ et $0{,}602$. 
L'apprentissage par renforcement avec rétroaction humaine (RLHF) 
dégrade significativement la calibration sur les familles Mistral-7B 
et Llama-3.1-8B (quatre tests, $p < 0{,}003$). Trois découvertes 
originales émergent de l'analyse approfondie : les prompts enrichis 
dégradent la calibration dans cinq modèles sur six ; cinq 
configurations présentent un effondrement de mode, prédisant un 
label unique dans plus de $97\%$ des cas ; sur les $238$ paires 
nécessitant une décision véritablement discriminante, les deux LLM 
médicaux testés (Meditron et BioMistral) atteignent une exactitude 
inférieure ou égale à $0{,}4\%$. La recalibration isotonic post-hoc 
ramène onze des douze configurations sous le seuil de calibration 
mais ne peut pas restaurer le raisonnement des modèles effondrés.

Les métriques agrégées de calibration et d'exactitude peuvent 
masquer un échec complet des LLM médicaux open source sur les cas 
cliniquement décisifs. Un protocole d'évaluation responsable doit 
combiner calibration, diversité prédictive et analyse stratifiée par 
difficulté de cas.

\vspace{0.5cm}
\noindent\textbf{Mots-clés :} grands modèles de langage, calibration, 
appariement d'essais cliniques, RLHF, IA médicale, mode collapse, 
expected calibration error.

\selectlanguage{english}
\cleardoublepage